% concatenated from Sendov_Scalar_Optimization_Proof.tex
\documentclass[11pt]{amsart}

\usepackage[T1]{fontenc}
\usepackage{lmodern}
\usepackage{microtype}
\usepackage{geometry}
\geometry{margin=1in}
\usepackage{amsmath,amssymb,amsthm,mathtools}
\usepackage{booktabs,array,longtable}
\usepackage{pdflscape}
\usepackage{xcolor}
\usepackage{enumitem}
\usepackage{hyperref}
\hypersetup{
  colorlinks=true,
  linkcolor=blue!55!black,
  citecolor=blue!55!black,
  urlcolor=blue!65!black,
  pdftitle={A Scalar Optimization Proof of Sendov's Conjecture with Reduced Computer Assistance},
  pdfauthor={Senjian An},
  pdfkeywords={Sendov's conjecture, polynomial critical points, scalar optimization, convexity, log-concavity, interval certificate}
}
\usepackage[nameinlink,noabbrev]{cleveref}
\usepackage{url}

% AMS classes capitalize the title and author by default.  The following
% local title-page definitions retain ordinary title case and place the
% affiliation and email immediately below the author's name.
\makeatletter
\def\@settitle{%
  \begin{center}
    \baselineskip=15pt\relax
    {\Large\bfseries\@title\par}
  \end{center}}
\def\@setauthors{%
  \begingroup
  \trivlist
  \centering
  \item\relax
  {\large\authors\par}
  \vspace{0.4em}
  {\normalfont\small
   School of Electrical Engineering, Computing and Mathematical Sciences,\par
   Curtin University, Perth, Australia\par
   \vspace{0.15em}
   \href{mailto:s.an@curtin.edu.au}{\texttt{s.an@curtin.edu.au}}\par}
  \endtrivlist
  \endgroup}
\makeatother

\newtheorem{theorem}{Theorem}[section]
\newtheorem{proposition}[theorem]{Proposition}
\newtheorem{lemma}[theorem]{Lemma}
\newtheorem{corollary}[theorem]{Corollary}
\newtheorem{claim}[theorem]{Claim}
\theoremstyle{definition}
\newtheorem{definition}[theorem]{Definition}
\newtheorem{remark}[theorem]{Remark}

\newcommand{\D}{\mathbb D}
\newcommand{\C}{\mathbb C}
\newcommand{\R}{\mathbb R}
\newcommand{\dd}{\,\mathrm d}
\newcommand{\e}{\mathrm e}
\newcommand{\etaStar}{\eta^{\ast}}
\newcommand{\betaHat}{\widehat\beta}
\newcommand{\Acal}{\mathcal A}

\title[Scalar optimization proof of Sendov's conjecture]{A Scalar Optimization Proof of Sendov's Conjecture with Reduced Computer Assistance}
\author{Senjian An}
\subjclass[2020]{Primary 30C10; Secondary 26C10, 26D15}
\keywords{Sendov's conjecture, polynomial critical points, scalar optimization, reduced computer assistance, convexity, log-concavity, exact certificate}
\date{}

\begin{document}

\begin{abstract}
Let $p$ be a complex polynomial of degree $n\ge2$ whose zeros lie in the closed unit disk.  Sendov's conjecture asserts that every zero of $p$ lies within distance one of a zero of $p'$.  Mazur's recent proof, and Tao's streamlined exposition of it, reduce a hypothetical counterexample to a scalar lower bound
\[
  1\le F(\eta,\alpha,n)
\]
together with two upper bounds for the endpoint parameter $\eta$ and the uniform restriction $0<\alpha\le17$.

We complete this scalar reduction by a two-stage optimization argument.  First, $F$ is nondecreasing in $\eta$, so $\eta$ may be replaced by a piecewise polar envelope $\eta^{\ast}(\alpha)$.  Writing $s=(n-1)/2$ converts the degree to a half-integer variable and gives
\[
  F(\eta^{\ast}(\alpha),\alpha,n)
  =E(\alpha,s)+\kappa(\alpha,s)
   \int_0^1 t^3\widehat\beta(t;\alpha,s)^{s-3/2}\dd t.
\]
On each half-unit $\alpha$-slab, $E$ and $\kappa$ decrease with $\alpha$, whereas $\widehat\beta$ increases.  Convexity in $t$ reduces all but the first mesh interval to eight explicit nodal terms.  The first interval satisfies a uniform bound $3/200$.  Each nodal upper term is a strictly log-concave function of the half-integer $s$, so its global discrete maximum is certified by two adjacent ratio evaluations.  A fixed finite certificate over the $34$ half-unit slabs gives
\[
  F(\eta^{\ast}(\alpha),\alpha,n)<0.966537<\frac{97}{100}<1,
\]
contradicting the scalar lower bound.  The main contribution is a substantial reduction of the computer assistance required after the Mazur--Tao scalar reduction: all continuous optimization and the unbounded degree parameter are handled analytically, leaving only a small fixed collection of explicit one-variable inequalities.  These are certified by exact rational bounds, with an independent interval cross-check.  Thus the computer-assisted part is confined to a transparent finite certificate rather than a large formal or numerical search.
\end{abstract}

\maketitle
\markboth{Senjian An}{Scalar optimization proof of Sendov's conjecture}

\section{Introduction}\label{sec:introduction}

Sendov's conjecture states that if every zero of a complex polynomial lies in the closed unit disk, then each zero has a critical point within unit distance.  More precisely, if $p\in\C[z]$ has degree $n\ge2$, all zeros of $p$ lie in $\overline\D$, and $a$ is a zero of $p$, then there is a zero $w$ of $p'$ with $|w-a|\le1$.

The conjecture was proved in successively larger low-degree and asymptotic regimes; see, among others, Meir--Sharma~\cite{MeirSharma1969}, Brown--Xiang~\cite{BrownXiang1999}, D\'egot~\cite{Degot2014}, and Tao~\cite{Tao2022}.  In 2026 Mazur presented a computer-assisted all-degree proof~\cite{Mazur2026}.  Tao subsequently gave a streamlined exposition and a separate Lean formalization of the underlying argument~\cite{TaoDigestion2026,TaoLean2026}.  The decisive feature of the Mazur--Tao approach is that a counterexample is reduced to a small scalar optimization problem.

The proof boundary is as follows.  The polar and origin inequalities in Section~\ref{sec:reduction}, together with the bound $0<\alpha\le17$, are imported from the Mazur--Tao argument.  The contribution of the present paper begins with the explicit scalar function $F(\eta,\alpha,n)$ and gives a new optimization argument designed specifically to reduce the amount of computer assistance needed to complete the proof.  The continuous parameter reductions, monotonicity arguments, convexity estimates, and the unbounded degree optimization are all handled analytically.  Only a fixed finite family of elementary one-variable inequalities remains for exact certification.

The new completion has four reusable ingredients:
\begin{enumerate}[label=(\roman*),leftmargin=2em]
\item monotonicity of the scalar objective in the endpoint parameter $\eta$;
\item a reparameterization $s=(n-1)/2$ that absorbs the degree into one half-integer variable;
\item monotonicity in $\alpha$ and convexity in the integration variable $t$;
\item strict log-concavity in $s$, reducing every remaining maximum to two adjacent ratio checks and one function evaluation.
\end{enumerate}
No stationary-point analysis in a two-dimensional parameter space is required.  The analytic reductions are proved in the text, and strict log-concavity removes the need for an unbounded search over the degree.  The remaining computer assistance is therefore deliberately small: it consists only of directed evaluations of a fixed finite list of explicit one-variable functions.  The corresponding exact certificate and independent checking scripts are included with the source.  In this sense, the novelty of the present proof is not the introduction of further computation, but the reduction of the computational burden to a compact and auditable certificate.

\paragraph{Main contribution.}
Relative to the existing computer- and Lean-assisted proofs, the present argument shifts the proof burden from formal or large-scale computational verification to elementary scalar analysis.  After the Mazur--Tao reduction, every continuous optimization is resolved analytically, and the only residual computer-assisted step is the exact verification of a small fixed certificate.  This reduced-certificate structure makes the proof easier to inspect and reproduce while confining all computer assistance to a fixed exact certificate.

\begin{theorem}[Sendov's conjecture]\label{thm:sendov}
Let $p\in\C[z]$ have degree $n\ge2$, and suppose every zero of $p$ lies in $\overline\D$.  For every zero $a$ of $p$, there exists a zero $w$ of $p'$ such that $|w-a|\le1$.
\end{theorem}

The classical low-degree arguments cover $n\le5$.  We therefore work with the Mazur--Tao scalar reduction for $n\ge6$.

\section{The Mazur--Tao scalar reduction}\label{sec:reduction}

We state the precise input from Tao's exposition~\cite{TaoDigestion2026}.  Suppose, for contradiction, that a counterexample exists.  The central and boundary cases $a=0$ and $|a|=1$ are already handled in the Mazur--Tao framework; after rotation we may therefore assume that the distinguished zero is real, with $0<a<1$.  If the critical points are $w_1,\ldots,w_{n-1}$, failure of Sendov's conclusion gives $|w_j-a|>1$; hence, with multiplicity,
\begin{equation}\label{eq:q-param}
 w_j=a-\frac1{q_j},\qquad 0<|q_j|<1,
 \qquad 1\le j\le n-1.
\end{equation}
Define
\begin{equation}\label{eq:x-def}
 x+i\xi=\frac1{n-1}\sum_{j=1}^{n-1}q_j
\end{equation}
and
\begin{equation}\label{eq:scalar-params}
 \alpha=\frac{n-1}{2}(1-a^2),\qquad
 \beta(t)=1-2atx+a^2t^2,\qquad
 \eta=\beta(1).
\end{equation}
The disk is convex, so $|x+i\xi|<1$.

\begin{proposition}[Simplified polar inequality; Tao, Proposition 10(iii)]\label{prop:polar}
For the quantities in \eqref{eq:x-def}--\eqref{eq:scalar-params},
\begin{equation}\label{eq:polar}
 0<\eta<\min\left\{
 \frac{\alpha}{\alpha+3},\,
 1-\frac{\log\alpha}{\alpha}
 \right\}.
\end{equation}
Moreover,
\[
 x>\frac a2>0.
\]
\end{proposition}

This is the ``Simplified polar inequality'' in Proposition 10(iii) of Tao's blog~\cite{TaoDigestion2026}.  The last assertion is equivalent to
\[
 1-\eta=2a\left(x-\frac a2\right)>0.
\]

From \eqref{eq:scalar-params},
\[
 a^2=1-\frac{2\alpha}{n-1},\qquad
 ax=1-\frac{\alpha}{n-1}-\frac\eta2.
\]
Thus
\begin{equation}\label{eq:beta-eta-alpha-n}
 \beta(t)=\beta_{\eta,\alpha,n}(t)
 :=1-2\left(1-\frac{\alpha}{n-1}-\frac\eta2\right)t
 +\left(1-\frac{2\alpha}{n-1}\right)t^2.
\end{equation}

\begin{proposition}[Alpha bound and origin inequality; Tao, Proposition 11(ii),(iii)]\label{prop:origin}
For $n\ge5$,
\[
 0<\alpha\le17,
\]
and
\begin{align}
1\le{}&
 \frac{\eta}{2\alpha(1-\eta)}
 +\frac{\eta}{4\alpha}
 +\frac1{2(n-1)}
 +\frac{\eta}{4\alpha(n-1)}\notag\\
&+\frac{\left(1-\frac{2\alpha}{n-1}\right)^2
 n(n-1)(n-2)\eta}{4\alpha}
 \int_0^1 t^3
 \beta_{\eta,\alpha,n}(t)^{(n-4)/2}\dd t\notag\\
=:{}&F(\eta,\alpha,n).
\label{eq:origin}
\end{align}
\end{proposition}

The bound on $\alpha$ is Proposition 11(ii), and the inequality is Proposition 11(iii), of Tao's exposition~\cite{TaoDigestion2026}.  Everything below concerns only the explicit scalar function in \eqref{eq:origin}.

\section{The scalar optimization problem}\label{sec:optimization}

Put
\[
 A=1-\frac{2\alpha}{n-1}=a^2,
 \qquad
 c=1-\frac{\alpha}{n-1}-\frac\eta2=ax.
\]
The original geometry gives
\[
 0<A<1,
 \qquad 0<c\le\sqrt A,
\]
and therefore
\begin{equation}\label{eq:beta-positive}
 \beta_{\eta,\alpha,n}(t)
 =(1-ct)^2+t^2(A-c^2)>0,
 \qquad 0\le t\le1.
\end{equation}
Every counterexample gives $n\ge6$ and an admissible pair $(\alpha,\eta)$ satisfying \eqref{eq:polar}, $0<\alpha\le17$, and $1\le F(\eta,\alpha,n)$.

The first observation allows the endpoint variable to be eliminated.

\begin{lemma}[Monotonicity in $\eta$]\label{lem:Feta}
For fixed admissible $\alpha$ and $n$,
\[
 \frac{\partial}{\partial\eta}F(\eta,\alpha,n)\ge0.
\]
\end{lemma}

The proof is given in Appendix~\ref{app:eta}.  The key identity is
\[
 \frac{\partial}{\partial\eta}\beta_{\eta,\alpha,n}(t)=t\ge0,
\]
and the appendix differentiates the complete integral factor rather than only the quadratic.

Define the polar envelope
\begin{equation}\label{eq:eta-star}
 \etaStar(\alpha)=
 \begin{cases}
 \dfrac{\alpha}{\alpha+3},&0<\alpha\le11,\\[2mm]
 1-\dfrac{\log\alpha}{\alpha},&11<\alpha\le17.
 \end{cases}
\end{equation}
By Proposition~\ref{prop:polar}, $0<\eta<\etaStar(\alpha)<1$.

\begin{lemma}[Preservation of feasibility]\label{lem:feasibility}
Let
\[
 c_0=1-\frac{\alpha}{n-1}-\frac\eta2,
 \qquad
 c^{\ast}=1-\frac{\alpha}{n-1}-\frac{\etaStar(\alpha)}2.
\]
Then
\[
 0<c^{\ast}\le c_0\le\sqrt A,
 \qquad (c^{\ast})^2\le A.
\]
\end{lemma}

\begin{proof}
Since $\eta\le\etaStar$, one has $c^{\ast}\le c_0\le\sqrt A$.  Moreover,
\[
 c^{\ast}
 =\frac{1+A-\etaStar}{2}
 >\frac A2>0
\]
because $\etaStar<1$.  Since $c^{\ast}>0$ and $c^{\ast}\le\sqrt A$, squaring gives $(c^{\ast})^2\le A$.  This proves the claim.
\end{proof}

In particular, $0<c^{\ast}\le\sqrt A<1$.  Hence for $0\le t\le1$,
\[
 1-c^{\ast}t\ge1-c^{\ast}>0,
\]
and therefore
\[
 \beta_{\etaStar(\alpha),\alpha,n}(t)
 =(1-c^{\ast}t)^2+t^2\bigl(A-(c^{\ast})^2\bigr)>0.
\]
This remains strict even in the limiting case $A=(c^{\ast})^2$.

Lemmas~\ref{lem:Feta} and~\ref{lem:feasibility} justify the replacement $\eta\mapsto\etaStar(\alpha)$.  We now absorb the degree into
\begin{equation}\label{eq:s-def}
 s=\frac{n-1}{2}
 \in\left\{\frac52,3,\frac72,\ldots\right\},
 \qquad
 p=s-\frac32,
 \qquad
 \alpha<s.
\end{equation}
The last inequality is equivalent to $a^2=1-\alpha/s>0$.

After the substitution,
\begin{equation}\label{eq:beta-hat}
 \betaHat(t;\alpha,s)
 =(1-t)^2+\etaStar(\alpha)t
 +\frac{\alpha}{s}t(1-t),
\end{equation}
\begin{equation}\label{eq:E}
 E(\alpha,s)=
 \frac{\etaStar}{2\alpha(1-\etaStar)}
 +\frac{\etaStar}{4\alpha}
 +\frac1{4s}
 +\frac{\etaStar}{8\alpha s},
\end{equation}
and
\begin{equation}\label{eq:kappa}
 \kappa(\alpha,s)=
 \frac{(4s^2-1)(s-\alpha)^2}{2\alpha s}\,\etaStar(\alpha).
\end{equation}
Consequently
\begin{equation}\label{eq:Fstar}
 F(\etaStar(\alpha),\alpha,n)
 =E(\alpha,s)+\kappa(\alpha,s)
 \int_0^1 t^3\betaHat(t;\alpha,s)^p\dd t.
\end{equation}
The proof of Sendov's conjecture is therefore reduced to proving that the right-hand side of \eqref{eq:Fstar} is uniformly smaller than one on
\[
 0<\alpha\le17,
 \quad
 s\in\tfrac12\mathbb Z,
 \quad
 s\ge\tfrac52,
 \quad
 \alpha<s.
\]

\section{Monotonicity on half-unit alpha slabs}\label{sec:slabs}

Partition $(0,17]$ into the $34$ half-unit slabs
\begin{equation}\label{eq:slabs}
 I_m=\left(\frac m2,\frac{m+1}{2}\right],
 \qquad m=0,\ldots,33.
\end{equation}
Write $I_m=(a,b]$.  Since $b$ and $s$ are half-integers and $\alpha<s$, every feasible point in the slab satisfies
\begin{equation}\label{eq:s-ge-b}
 s\ge b.
\end{equation}
The interval $(10.5,11]$ uses the rational branch of \eqref{eq:eta-star}, whereas $(11,11.5]$ uses the logarithmic branch; no slab crosses the change of formula.

\begin{lemma}[Monotonicity of the three factors]\label{lem:slab-monotone}
For fixed $s$ and $t\in[0,1]$, on either branch of \eqref{eq:eta-star},
\[
 \partial_\alpha\betaHat(t;\alpha,s)\ge0,
 \qquad
 \partial_\alpha E(\alpha,s)<0,
 \qquad
 \partial_\alpha\kappa(\alpha,s)<0.
\]
The first inequality is strict whenever $t>0$.
\end{lemma}

\begin{proof}
On $0<\alpha\le11$,
\[
 \etaStar=\frac{\alpha}{\alpha+3},
\]
so
\[
 \partial_\alpha\betaHat
 =\frac{3t}{(\alpha+3)^2}+\frac{t(1-t)}s\ge0.
\]
Equality occurs only at $t=0$.
\[
 E=\frac16+\frac1{4(\alpha+3)}
 +\frac1{4s}+\frac1{8s(\alpha+3)},
\]
and
\[
 \kappa
 =\frac{(4s^2-1)(s-\alpha)^2}{2s(\alpha+3)}.
\]
The last two expressions are strictly decreasing in $\alpha$.

On $11<\alpha\le17$, write $L=\log\alpha$.  Then
\[
 \etaStar=1-\frac L\alpha,
 \qquad
 \partial_\alpha\betaHat
 =\frac{t(L-1)}{\alpha^2}+\frac{t(1-t)}s\ge0.
\]
Again equality occurs only at $t=0$, because $L=\log\alpha>1$ for $\alpha>11$.
Furthermore
\[
 \frac{\etaStar}{\alpha}
 =\frac{\alpha-L}{\alpha^2},
 \qquad
 \left(\frac{\etaStar}{\alpha}\right)'
 =-\frac{\alpha+1-2L}{\alpha^3}<0.
\]
Indeed, $L<3$ and $\alpha>11$, so
\[
 \alpha+1-2L>12-6=6>0.
\]
Both $(s-\alpha)^2$ and $\etaStar/\alpha$ are positive and decreasing, so $\kappa$ decreases.

Finally,
\[
 E=
 \frac1{2L}-\frac1{4\alpha}-\frac{L}{4\alpha^2}
 +\frac1{4s}+\frac1{8\alpha s}-\frac{L}{8\alpha^2s}.
\]
The derivative of the first three terms is
\[
 \frac1{4\alpha^2}
 \left(1+\frac{2L-1}{\alpha}-\frac{2\alpha}{L^2}\right).
\]
Since $L<3$ and $\alpha>11$,
\[
 1+\frac{2L-1}{\alpha}
 <1+\frac5{11}=\frac{16}{11},
 \qquad
 \frac{2\alpha}{L^2}>\frac{22}{9},
\]
and $16/11<22/9$; hence this derivative is strictly negative.  The derivative of the last two $\alpha$-dependent terms is
\[
 -\frac{\alpha-2L+1}{8s\alpha^3}<0,
\]
because $\alpha-2L+1>6$.
\end{proof}

For a slab $(a,b]$, let $E_a(s)$ and $\kappa_a(s)$ denote the right limits at $a$ on the branch belonging to the slab.  At $a=0$ these limits are
\begin{equation}\label{eq:a0-limits}
 E_0(s)=\frac14+\frac7{24s},
 \qquad
 \kappa_0(s)=\frac{(4s^2-1)s}{6}.
\end{equation}
Lemma~\ref{lem:slab-monotone} gives
\begin{equation}\label{eq:slab-majorization}
 E(\alpha,s)\le E_a(s),
 \qquad
 \kappa(\alpha,s)\le\kappa_a(s),
 \qquad
 \betaHat(t;\alpha,s)\le\betaHat(t;b,s).
\end{equation}

\section{Decomposition of the integral}\label{sec:decomposition}

Let
\[
 t_j=\frac j8,
 \qquad j=0,\ldots,8.
\]
Because $\alpha<s$,
\[
 \frac{\partial^2\betaHat}{\partial t^2}
 =2\left(1-\frac\alpha s\right)>0.
\]
Moreover $p=s-3/2\ge1$, so
\begin{equation}\label{eq:power-convex}
 \frac{\dd^2}{\dd t^2}\betaHat(t)^p
 =p\betaHat^{p-2}
 \left((p-1)(\betaHat')^2+\betaHat\betaHat''\right)\ge0.
\end{equation}
Thus $t\mapsto\betaHat(t)^p$ is convex.

\subsection{The first interval}

The interval $[0,1/8]$ is treated by retaining the full power of the chord of $\sqrt{\betaHat}$.  Since
\[
 \frac{\dd^2}{\dd t^2}\sqrt{\betaHat(t)}
 =\frac{A-(c^{\ast})^2}{\betaHat(t)^{3/2}}\ge0,
\]
its graph lies below the chord joining $t=0$ and $t=1/8$.  Put
\[
 r_1=\sqrt{\betaHat(1/8)},
 \qquad
 \ell_0(t)=1-8(1-r_1)t,
 \qquad 0\le t\le\frac18.
\]
Then $\sqrt{\betaHat(t)}\le\ell_0(t)$ and, because $2p=2s-3$,
\[
 \kappa(\alpha,s)\int_0^{1/8}t^3\betaHat(t)^p\dd t
 \le T_0(\alpha,s),
\]
where the first-interval chord majorant is defined by
\begin{equation}\label{eq:T0-def}
 T_0(\alpha,s)
 :=\kappa(\alpha,s)\int_0^{1/8}t^3\ell_0(t)^{2s-3}\dd t.
\end{equation}

\begin{lemma}[First-interval estimate]\label{lem:T0}
For every admissible $\alpha$ and $s$,
\[
 T_0(\alpha,s)<\frac3{200}.
\]
\end{lemma}

A complete one-variable proof is given in Appendix~\ref{app:T0}.  It reduces $T_0$ to a Laplace integral and uses only rational inequalities and finite exponential series.

\subsection{The remaining seven intervals}

On $[t_i,t_{i+1}]$, $i=1,\ldots,7$, convexity \eqref{eq:power-convex} gives
\[
 \betaHat(t)^p
 \le
 8(t_{i+1}-t)\betaHat(t_i)^p
 +8(t-t_i)\betaHat(t_{i+1})^p.
\]
Multiplying by $t^3$ and integrating, then collecting adjacent endpoint contributions, yields
\begin{equation}\label{eq:nodal-decomp}
 \int_{1/8}^{1}t^3\betaHat(t)^p\dd t
 \le\sum_{j=1}^{8}c_j\betaHat(t_j)^p,
\end{equation}
where
\begin{equation}\label{eq:weights}
(c_1,\ldots,c_8)=
\left(
\frac{13}{40960},\frac9{4096},\frac{57}{8192},\frac{33}{2048},
\frac{255}{8192},\frac{219}{4096},\frac{693}{8192},
\frac{4519}{81920}
\right).
\end{equation}
The elementary integrations producing these weights are recorded in Appendix~\ref{app:weights}.

Fix a slab $(a,b]$.  Define
\begin{equation}\label{eq:Uj}
 U_j^{a,b}(s)=
 c_j\kappa_a(s)
 \left((1-t_j)^2+\etaStar(b)t_j
 +\frac{b}{s}t_j(1-t_j)\right)^{s-3/2}.
\end{equation}
Combining \eqref{eq:slab-majorization}, Lemma~\ref{lem:T0}, and \eqref{eq:nodal-decomp} gives the master slab bound
\begin{equation}\label{eq:master-slab}
 F(\etaStar(\alpha),\alpha,n)
 <E_a(s)+\frac3{200}+\sum_{j=1}^{8}U_j^{a,b}(s).
\end{equation}

\section{One-dimensional maxima and the finite certificate}\label{sec:certificate}

For fixed $a,b,j$, put
\[
 B_0=(1-t_j)^2+\etaStar(b)t_j,
 \qquad
 B_1=bt_j(1-t_j).
\]
Apart from the positive constant $c_j$, the term in \eqref{eq:Uj} is
\begin{equation}\label{eq:P}
 P(s)=\kappa_a(s)\left(B_0+\frac{B_1}{s}\right)^{s-3/2}.
\end{equation}
The admissible values form the half-integer lattice
\begin{equation}\label{eq:lattice}
 \Lambda_b
 =\left\{s_0+\frac{k}{2}:k=0,1,2,\ldots\right\},
 \qquad
 s_0=\max\left\{\frac52,b\right\}.
\end{equation}
Here $b$ is itself a half-integer.  The condition $s\ge b$ follows from
$\alpha\in(a,b]$, $\alpha<s$, and the half-integrality of $s$.

\begin{lemma}[Strict log-concavity]\label{lem:logconcavity}
For $s>\max\{a,1/2\}$, the function $P$ in \eqref{eq:P} is strictly log-concave.
\end{lemma}

\begin{proof}
Let
\[
 \phi(s)=\log P(s).
\]
The positive constants independent of $s$ disappear after differentiation.  Since
\[
 \kappa_a(s)=C_a\,
 \frac{(s-\frac12)(s+\frac12)(s-a)^2}{s}
\]
for a positive constant $C_a$, one has
\[
 \log\kappa_a(s)
 =\log C_a+
 \log\left(s-\frac12\right)
 +\log\left(s+\frac12\right)
 +2\log(s-a)-\log s.
\]
Consequently,
\begin{equation}\label{eq:logkappa2}
 (\log\kappa_a)''
 =-\frac1{(s-\frac12)^2}
  -\frac1{(s+\frac12)^2}
  -\frac2{(s-a)^2}
  +\frac1{s^2}<0.
\end{equation}
Indeed, $s-1/2<s$ gives
$-1/(s-1/2)^2+1/s^2<0$, and the remaining two terms are nonpositive, with $-2/(s-a)^2<0$.

It remains to treat the power factor.  Set
\[
 q(s)=B_0+\frac{B_1}{s}=\frac{B_0s+B_1}{s},
 \qquad
 h(s)=\left(s-\frac32\right)\log q(s).
\]
Because
\[
 \frac{q'(s)}{q(s)}
 =-\frac{B_1}{s(B_0s+B_1)},
\]
we obtain
\[
 h'(s)
 =\log q(s)
 -\frac{(s-\frac32)B_1}{s(B_0s+B_1)}.
\]
Differentiating once more and putting the terms over the common denominator
$2s^2(B_0s+B_1)^2$ gives
\begin{equation}\label{eq:powerlog2}
 h''(s)
 =-\frac{B_1(6B_0s+2B_1s+3B_1)}
 {2s^2(B_0s+B_1)^2}\le0.
\end{equation}
All factors in the denominator are positive.  Also $B_0>0$ and $B_1\ge0$; when $t_j=1$, one has $B_1=0$ and hence $h''=0$.

Finally,
\[
 \phi''(s)=(\log\kappa_a)''(s)+h''(s)<0
\]
by \eqref{eq:logkappa2} and \eqref{eq:powerlog2}.  Thus $P$ is strictly log-concave.
\end{proof}

For any $s\in\Lambda_b$, define the adjacent ratio
\begin{equation}\label{eq:ratio-def}
 R(s)=\frac{P(s+1/2)}{P(s)}
 =\frac{\kappa_a(s+1/2)}{\kappa_a(s)}
 \frac{\left(B_0+B_1/(s+1/2)\right)^{s-1}}
      {\left(B_0+B_1/s\right)^{s-3/2}}.
\end{equation}
The ratio is defined at every lattice point because $s+1/2\in\Lambda_b$.
If $\phi=\log P$, then
\[
 \frac{\mathrm d}{\mathrm ds}\log R(s)
 =\phi'(s+1/2)-\phi'(s)<0,
\]
since strict concavity makes $\phi'$ strictly decreasing.  Therefore $R(s)$ is strictly decreasing on the whole continuous domain, and in particular on the half-integer lattice.

\begin{corollary}[Two-ratio maximum certificate]\label{cor:ratio}
Let $\Lambda_b$ and $s_0$ be as in \eqref{eq:lattice}, and let $s_\ast\in\Lambda_b$.
\begin{enumerate}[label=(\alph*),leftmargin=2.2em]
\item If $s_\ast>s_0$ and
\[
 R(s_\ast-1/2)\ge1,
 \qquad
 R(s_\ast)\le1,
\]
then $P(s_\ast)$ is a global maximum of $P$ on $\Lambda_b$.
\item If $s_\ast=s_0$, there is no admissible lattice point to its left.  In this boundary case the single condition
\[
 R(s_0)\le1
\]
is sufficient to conclude that $P(s_0)$ is a global maximum on $\Lambda_b$.
\end{enumerate}
The maximum need not be unique when one of the displayed inequalities is an equality.
\end{corollary}

\begin{proof}
First suppose $s_\ast>s_0$.  Let $r\in\Lambda_b$ with $r<s_\ast$.  Every ratio occurring while moving from $r$ to $s_\ast$ has its argument at most $s_\ast-1/2$.  Since $R$ is decreasing,
\[
 R(u)\ge R(s_\ast-1/2)\ge1
 \qquad
 (r\le u\le s_\ast-1/2,
 \ u\in\Lambda_b).
\]
Thus
\[
 P(u+1/2)=R(u)P(u)\ge P(u)
\]
at every step, and iteration gives $P(r)\le P(s_\ast)$.

Now let $r>s_\ast$.  Every ratio used while moving from $s_\ast$ to $r$ has argument at least $s_\ast$.  Therefore
\[
 R(u)\le R(s_\ast)\le1
 \qquad
 (s_\ast\le u\le r-1/2,
 \ u\in\Lambda_b),
\]
and hence $P(u+1/2)\le P(u)$ at every step.  Iteration again gives $P(r)\le P(s_\ast)$.  This proves part (a).

If $s_\ast=s_0$, the left-hand argument is unnecessary because no point of $\Lambda_b$ lies below $s_0$.  Since $R$ is decreasing, $R(u)\le R(s_0)\le1$ for every $u\in\Lambda_b$ with $u\ge s_0$.  Hence the sequence is nonincreasing from its first term onward, proving part (b).
\end{proof}

\begin{remark}[Boundary case]\label{rem:left-endpoint}
If $s_\ast=s_0$, then $s_0-1/2\notin\Lambda_b$, so there is no predecessor ratio to verify.  The condition $R(s_0)\le1$, together with the decrease of $R$, implies that $P$ is nonincreasing on the entire admissible lattice beginning at $s_0$.
\end{remark}

The elementary term $E_a(s)$ has the form $C_a+D_a/s$ with $D_a>0$, and is therefore maximized at the smallest admissible half-integer $s=\max\{5/2,b\}$.

Appendix~\ref{app:full-certificate} gives, for every slab and each $U_j$, the maximizing half-integer and an upward-rounded interval upper bound.  The archived file \path{Scalar_Certificate_HalfUnit.csv} is part of the finite certificate: in addition to the values printed in the PDF tables, it records a directed lower bound for $R(s_\ast-1/2)$ whenever $s_\ast>s_0$, and a directed upper bound for $R(s_\ast)$ in every case.  The fixed-candidate checker \path{verify_scalar_certificate_exact.py} reads these proposed maximizers and verifies the ratio signs, upper values, and row totals using exact rational arithmetic, rational Taylor enclosures for logarithms, and rational square-root brackets.  It verifies the stated maximizing degree without performing any search.

A representative check occurs on the worst slab $(0,1/2]$ for $U_1$.  Here the table lists $s_\ast=12.5$, and directed interval evaluation gives
\[
 R(12)>1.0009543562>1,
 \qquad
 R(12.5)<0.9961149406<1,
\]
while
\[
 U_1^{0,1/2}(12.5)<0.029954.
\]
Corollary~\ref{cor:ratio} proves that this is the global half-integer maximum.  Every other entry is certified in exactly the same way.

For convenience, the row totals are summarized below.  Each total includes the upward-rounded maximum of $E_a$, the eight upward-rounded nodal maxima, and $3/200$.

\begin{center}
\small
\begin{tabular}{cc@{\qquad\qquad}cc}
\toprule
$\alpha$ slab & upper total & $\alpha$ slab & upper total\\
\midrule
$(0,0.5]$ & 0.966537 & $(8.5,9]$ & 0.868742 \\
$(0.5,1]$ & 0.909810 & $(9,9.5]$ & 0.881991 \\
$(1,1.5]$ & 0.858561 & $(9.5,10]$ & 0.895303 \\
$(1.5,2]$ & 0.833062 & $(10,10.5]$ & 0.909358 \\
$(2,2.5]$ & 0.823678 & $(10.5,11]$ & 0.923668 \\
$(2.5,3]$ & 0.802545 & $(11,11.5]$ & 0.870814 \\
$(3,3.5]$ & 0.791615 & $(11.5,12]$ & 0.861055 \\
$(3.5,4]$ & 0.788032 & $(12,12.5]$ & 0.852305 \\
$(4,4.5]$ & 0.787550 & $(12.5,13]$ & 0.844273 \\
$(4.5,5]$ & 0.790978 & $(13,13.5]$ & 0.836964 \\
$(5,5.5]$ & 0.796623 & $(13.5,14]$ & 0.830001 \\
$(5.5,6]$ & 0.803976 & $(14,14.5]$ & 0.823685 \\
$(6,6.5]$ & 0.811866 & $(14.5,15]$ & 0.817881 \\
$(6.5,7]$ & 0.822444 & $(15,15.5]$ & 0.812479 \\
$(7,7.5]$ & 0.833053 & $(15.5,16]$ & 0.807370 \\
$(7.5,8]$ & 0.844389 & $(16,16.5]$ & 0.802504 \\
$(8,8.5]$ & 0.856239 & $(16.5,17]$ & 0.797954 \\
\bottomrule
\end{tabular}
\end{center}

The largest row is the first:
\begin{equation}\label{eq:global-certificate}
 E_0+\frac3{200}+\sum_{j=1}^8U_j^{0,1/2}
 <0.966537<\frac{97}{100}.
\end{equation}

\section{Completion of the proof}\label{sec:completion}

Suppose a counterexample exists in degree $n\ge6$.  Proposition~\ref{prop:origin} gives
\[
 1\le F(\eta,\alpha,n).
\]
By Lemma~\ref{lem:Feta} and the polar bound,
\[
 F(\eta,\alpha,n)
 \le F(\etaStar(\alpha),\alpha,n).
\]
Choose the half-unit slab $(a,b]$ containing $\alpha$.  The master bound \eqref{eq:master-slab}, the log-concavity certificate, and \eqref{eq:global-certificate} give
\[
 F(\etaStar(\alpha),\alpha,n)<\frac{97}{100}<1,
\]
a contradiction.  This proves Theorem~\ref{thm:sendov}.

\begin{remark}
The proof separates analysis from arithmetic.  The analytic part consists of monotonicity in $\eta$ and $\alpha$, convexity in $t$, and log-concavity in the half-integer $s$.  The arithmetic part consists of $34\times9$ explicit one-variable evaluations.  No numerical optimizer or unbounded degree search is used: strict log-concavity proves that each displayed half-integer is the global maximizer.
\end{remark}

\appendix

\section{Monotonicity in the endpoint parameter and feasibility}\label{app:eta}

Fix admissible $n$ and $\alpha$, and set $p=(n-4)/2>0$.  From \eqref{eq:beta-eta-alpha-n},
\[
 \frac{\partial}{\partial\eta}\beta_{\eta,\alpha,n}(t)=t\ge0.
\]
The elementary terms satisfy
\[
 \frac{\partial}{\partial\eta}
 \frac{\eta}{2\alpha(1-\eta)}
 =\frac1{2\alpha(1-\eta)^2}>0,
\]
\[
 \frac{\partial}{\partial\eta}\frac\eta{4\alpha}
 =\frac1{4\alpha}>0,
 \qquad
 \frac{\partial}{\partial\eta}\frac\eta{4\alpha(n-1)}
 =\frac1{4\alpha(n-1)}>0.
\]
For the integral contribution, define
\[
 I(\eta)=\eta\int_0^1t^3\beta_{\eta,\alpha,n}(t)^p\dd t.
\]
Differentiation under the integral sign gives
\[
 I'(\eta)=
 \int_0^1t^3\beta^{p-1}
 \bigl(\beta+\eta pt\bigr)\dd t\ge0.
\]
The remaining prefactor is positive, proving Lemma~\ref{lem:Feta}.

For completeness, the feasibility statement follows from
\[
 c^{\ast}=c_0-\frac{\etaStar-\eta}{2}\le c_0\le\sqrt A
\]
and
\[
 c^{\ast}=\frac{1+A-\etaStar}{2}>\frac A2>0.
\]
Thus the replacement $\eta\mapsto\etaStar$ both increases the objective and preserves positivity of the quadratic.

\section{The first-interval estimate}\label{app:T0}

We give the details of Lemma~\ref{lem:T0}.  Put
\[
 k=n-4=2s-3,
 \qquad
 B=2(\alpha+3),
 \qquad
 w=n-1-2\alpha.
\]
For this estimate alone, the rational polar bound $\eta\le\alpha/(\alpha+3)$ may be used on the entire range $0<\alpha\le17$.  Indeed, the prefactor $\kappa$ is proportional to $\eta$, the nodal value $r_1=\sqrt{\betaHat(1/8)}$ increases with $\eta$, and the first-chord integral increases with $r_1$.  Thus replacing $\eta$ by the rational upper bound can only increase $T_0$.
Let $r_1=\sqrt{\betaHat(1/8)}$.  Substitution gives
\begin{equation}\label{eq:T0-X}
 X:=1-r_1^2
 =\frac{7w}{64(k+3)}+\frac3{4B},
 \qquad 0<X<\frac{15}{64},
\end{equation}
and
\begin{equation}\label{eq:T0-P}
 \frac\kappa{4096}
 <P:=\frac{(k+3)w^2}{8192B}.
\end{equation}
Define
\[
 D(X)=\frac X2+\frac{X^2}{8}+\frac{X^3}{16}.
\]
Since
\[
 \left(1-\frac X2-\frac{X^2}{8}-\frac{X^3}{16}\right)^2-(1-X)
 =\frac{X^4(X^2+4X+20)}{256}>0,
\]
and both sides are nonnegative,
\[
 r_1=\sqrt{1-X}
 \le 1-D(X),
 \qquad
 1-r_1\ge D(X).
\]
By \eqref{eq:T0-def}, the chord is $\ell_0(t)=1-8(1-r_1)t$.  With the change of variables $u=8t$ and $k=2s-3=n-4$,
\begin{align*}
 T_0
 &=\frac{\kappa}{4096}
   \int_0^1u^3\bigl(1-(1-r_1)u\bigr)^k\dd u\\
 &\le\frac{\kappa}{4096}
   \int_0^1u^3\bigl(1-D(X)u\bigr)^k\dd u.
\end{align*}
For $0\le u\le1$, the factor $1-D(X)u$ is positive and $1-z\le\e^{-z}$ gives
\[
 \bigl(1-D(X)u\bigr)^k
 \le\e^{-kD(X)u}.
\]
Using $\kappa/4096<P$ from \eqref{eq:T0-P}, and setting
\[
 \lambda=kD(X),
 \qquad
 \Psi(\lambda)=\int_0^1u^3\e^{-\lambda u}\dd u,
\]
we obtain the explicit reduction
\begin{equation}\label{eq:T0-Laplace-reduction}
 T_0<P\Psi(\lambda).
\end{equation}

Write
\[
 a_1=\frac{7w}{64(k+3)},
 \qquad
 a_2=\frac3{4B},
 \qquad X=a_1+a_2.
\]
Weighted AM--GM gives $X^3\ge(27/4)a_1^2a_2$, and hence
\[
 P\le\frac{8(k+3)^3}{3969}X^3.
\]
With $S(X)=X^2+2X+8$, one has $D(X)=XS(X)/16$ and
\[
 \frac{P}{\lambda^3}
 \le\frac{32768}{3969}
 \left(\frac1{S(X)}+\frac{3X}{16\lambda}\right)^3.
\]
If $\lambda\ge L$, then
\[
 \frac1{S(X)}+\frac{3X}{16\lambda}
 \le h_L(X):=\frac1{S(X)}+\frac{3X}{16L}.
\]
For $0\le X\le15/64$ and $0<L\le23/5$,
\[
 h_L'(X)
 =-\frac{2(X+1)}{S(X)^2}+\frac{3}{16L}.
\]
The elementary bounds
\[
 2(X+1)\le\frac{79}{32},
 \qquad
 S(X)^2\ge64,
 \qquad
 \frac3{16L}\ge\frac{15}{368}
\]
give
\[
 h_L'(X)
 \ge-\frac{79}{2048}+\frac{15}{368}
 =\frac{103}{47104}>0.
\]
Thus $h_L$ is increasing on the whole permitted $X$-interval.  Since
\[
 S(15/64)=\frac{34913}{4096},
\]
we conclude that
\[
 h_L(X)
 \le h_L(15/64)
 =\frac{4096}{34913}+\frac{45}{1024L}.
\]
Consequently, whenever $\lambda\ge L$ and $0<L\le23/5$,
\begin{equation}\label{eq:T0-C}
 \frac{P}{\lambda^3}
 \le C(L):=\frac{32768}{3969}
 \left(\frac{4096}{34913}+\frac{45}{1024L}\right)^3.
\end{equation}

Set
\[
 g(\lambda)=\lambda^3\Psi(\lambda)
 =\frac{6-\e^{-\lambda}(\lambda^3+3\lambda^2+6\lambda+6)}{\lambda}.
\]
We record the two elementary properties used below:
\begin{equation}\label{eq:T0-gprops}
 g'(\lambda)>0\quad(0<\lambda\le23/5),
 \qquad
 g(\lambda)<\frac89\quad(\lambda>0).
\end{equation}
For the first, the numerator of $g'$ is
\[
 H(\lambda)=
 \e^{-\lambda}(\lambda^4+\lambda^3+3\lambda^2+6\lambda+6)-6,
\]
with
\[
 H(0)=0,
 \qquad
 H'(\lambda)=\e^{-\lambda}\lambda^3(3-\lambda).
\]
Thus $H$ increases on $(0,3)$ and decreases thereafter.  It is enough to check $H(23/5)>0$.  The finite exponential estimate
\[
 \e^{23/5}<
 \sum_{m=0}^{9}\frac{(23/5)^m}{m!}
 +\frac{(23/5)^{10}}{10!\,[1-(23/5)/11]}<100
\]
gives
\[
 \e^{-23/5}
 \left[
 \left(\frac{23}{5}\right)^4+
 \left(\frac{23}{5}\right)^3+
 3\left(\frac{23}{5}\right)^2+
 6\left(\frac{23}{5}\right)+6
 \right]
 >\frac1{100}\frac{401351}{625}>6.
\]
Hence $H>0$ on $(0,23/5]$.

For the second property in \eqref{eq:T0-gprops}, define
\[
 Q(\lambda)=
 \e^{-\lambda}(\lambda^3+3\lambda^2+6\lambda+6)
 -6+\frac89\lambda.
\]
Then $g<8/9$ is equivalent to $Q>0$, and
\[
 Q'(\lambda)=\frac89-\lambda^3\e^{-\lambda},
 \qquad
 Q''(\lambda)=\lambda^2\e^{-\lambda}(\lambda-3).
\]
The only possible positive minimum occurs at the larger zero $\rho>3$ of $Q'$.  The finite Taylor estimate
\[
 \e^{29/6}<
 \sum_{m=0}^{7}\frac{(29/6)^m}{m!}
 +\frac{(29/6)^8}{8!\,[1-(29/6)/9]}
 <\frac{24389}{192}
 =\frac98\left(\frac{29}{6}\right)^3
\]
shows $\rho>29/6$.  At $Q'(\rho)=0$,
\[
 Q(\rho)=\frac89
 \left(\rho+1+\frac3\rho+\frac6{\rho^2}+\frac6{\rho^3}\right)-6.
\]
The expression in parentheses is increasing for $\rho\ge29/6$, and at $29/6$ it is larger than $27/4$.  Therefore $Q(\rho)>0$, proving $g<8/9$.

For $0<\lambda<1/4$, the estimate $D(X)\ge X/2$ gives
\[
 \lambda=kD(X)\ge\frac{kX}{2},
 \qquad
 X\le\frac{2\lambda}{k}.
\]
Consequently,
\begin{align*}
 P
 &\le\frac{8(k+3)^3}{3969}X^3\\
 &\le\frac{64}{3969}
    \left(\frac{k+3}{k}\right)^3\lambda^3\\
 &\le\frac{1000}{3969}\lambda^3
 <\frac{125}{31752},
\end{align*}
where the penultimate inequality uses $k\ge2$ and $(k+3)/k\le5/2$.  Since $\Psi(\lambda)\le\int_0^1u^3\dd u=1/4$, \eqref{eq:T0-Laplace-reduction} yields
\[
 T_0<\frac{125}{127008}<\frac3{200}.
\]
For $\lambda\ge23/5$, \eqref{eq:T0-C} and $g<8/9$ give the same conclusion.  The compact interval $[1/4,23/5]$ is covered by the exact rational rows in Table~\ref{tab:T0}.  On a row $L\le\lambda\le U$, monotonicity of $g$ gives
\[
 T_0\le C(L)g(U)<C(L)G<\frac3{200}.
\]
To verify $g(U)<G$, write
\[
 A(U)=U^3+3U^2+6U+6.
\]
Since $6-GU>0$, the inequality $g(U)<G$ is equivalent to
\begin{equation}\label{eq:g-exp-equivalent}
 \e^U<\frac{A(U)}{6-GU}.
\end{equation}
For the truncation order $N$ displayed in Table~\ref{tab:T0}, positivity of the exponential series and the geometric bound on the tail give
\begin{equation}\label{eq:exp-rational-upper}
 \e^U
 <\sum_{m=0}^{N}\frac{U^m}{m!}
 +\frac{U^{N+1}}{(N+1)!\,[1-U/(N+2)]}.
\end{equation}
Every quantity on the right of \eqref{eq:exp-rational-upper} is rational.  Direct rational comparison verifies that it is smaller than the right-hand side of \eqref{eq:g-exp-equivalent}; a second exact rational comparison gives $C(L)G<3/200$.  The file \path{verify_T0_certificate.py} performs precisely these two checks and writes the exact margins to \path{T0_Certificate.csv}.  This proves Lemma~\ref{lem:T0}.

\begin{table}[ht]
\centering
\small
\caption{Finite exact-rational cover for the first-interval bound.  The final column is the truncation order in \eqref{eq:exp-rational-upper}.}\label{tab:T0}
\begin{tabular}{cccc}
\toprule
$L$&$U$&$G$&$N$\\
\midrule
$1/4$&$3/4$&$3/50$&4\\
$3/4$&$3/2$&$263/1000$&5\\
$3/2$&$9/4$&$509/1000$&6\\
$9/4$&$3$&$353/500$&7\\
$3$&$17/5$&$39/50$&8\\
$17/5$&$18/5$&$809/1000$&7\\
$18/5$&$19/5$&$104/125$&8\\
$19/5$&$39/10$&$421/500$&8\\
$39/10$&$4$&$17/20$&9\\
$4$&$41/10$&$429/500$&8\\
$41/10$&$21/5$&$108/125$&9\\
$21/5$&$43/10$&$87/100$&8\\
$43/10$&$22/5$&$437/500$&9\\
$22/5$&$89/20$&$2189/2500$&9\\
$89/20$&$23/5$&$22/25$&9\\
\bottomrule
\end{tabular}
\end{table}

\section{Chord weights}\label{app:weights}

For $i=1,\ldots,7$, set
\[
 L_i=\int_{t_i}^{t_{i+1}}t^3\,8(t_{i+1}-t)\dd t,
 \qquad
 R_i=\int_{t_i}^{t_{i+1}}t^3\,8(t-t_i)\dd t.
\]
Direct integration gives
\[
\begin{array}{c|cc}
 i&L_i&R_i\\ \hline
1&13/40960&49/81920\\
2&131/81920&97/40960\\
3&47/10240&499/81920\\
4&821/81920&1/80\\
5&763/40960&1829/81920\\
6&2551/81920&1487/40960\\
7&989/20480&4519/81920.
\end{array}
\]
Then $c_1=L_1$, $c_j=R_{j-1}+L_j$ for $2\le j\le7$, and $c_8=R_7$, yielding \eqref{eq:weights}.

\section{Full finite certificate}\label{app:full-certificate}

A cell $0.029954\,@\,12.5$ means that the term is strictly smaller than $0.029954$ for every admissible half-integer $s$, and that its global maximizing half-integer is $s=12.5$.  All displayed bounds are rounded upward to six decimal places.  The final column adds the displayed $E$ bound, the eight displayed nodal bounds, and $3/200$.

\begin{landscape}
\scriptsize
\setlength{\tabcolsep}{2.8pt}
\renewcommand{\arraystretch}{1.08}
\begin{center}
\textbf{Table D.1. Rational branch, $0<\alpha\le5.5$.}\par\medskip
\begin{tabular}{c|c|cccccccc|c}
\toprule
$\alpha$ slab & $E$ & $U_1$ & $U_2$ & $U_3$ & $U_4$ & $U_5$ & $U_6$ & $U_7$ & $U_8$ & total\\
\midrule
$(0,0.5]$ & 0.366667\,@\,2.5 & 0.029954\,@\,12.5 & 0.034952\,@\,6 & 0.045084\,@\,4 & 0.061696\,@\,3 & 0.086158\,@\,2.5 & 0.110753\,@\,2.5 & 0.137467\,@\,2.5 & 0.078806\,@\,2.5 & 0.966537 \\
$(0.5,1]$ & 0.352381\,@\,2.5 & 0.030590\,@\,13.5 & 0.035086\,@\,7 & 0.043651\,@\,4.5 & 0.056625\,@\,3.5 & 0.074461\,@\,3 & 0.097295\,@\,3 & 0.129068\,@\,2.5 & 0.075653\,@\,2.5 & 0.909810 \\
$(1,1.5]$ & 0.341667\,@\,2.5 & 0.031063\,@\,15 & 0.035307\,@\,8 & 0.043350\,@\,5.5 & 0.054830\,@\,4.5 & 0.069569\,@\,4 & 0.088414\,@\,3.5 & 0.113690\,@\,3.5 & 0.065671\,@\,3.5 & 0.858561 \\
$(1.5,2]$ & 0.333334\,@\,2.5 & 0.031415\,@\,16 & 0.035654\,@\,8.5 & 0.043149\,@\,6.5 & 0.053582\,@\,5 & 0.067034\,@\,4.5 & 0.084233\,@\,4.5 & 0.106897\,@\,4.5 & 0.062764\,@\,4.5 & 0.833062 \\
$(2,2.5]$ & 0.326667\,@\,2.5 & 0.031654\,@\,17 & 0.035941\,@\,9.5 & 0.043366\,@\,7 & 0.053550\,@\,6 & 0.066268\,@\,5.5 & 0.082552\,@\,5 & 0.105741\,@\,5 & 0.062939\,@\,5.5 & 0.823678 \\
$(2.5,3]$ & 0.303031\,@\,3 & 0.031803\,@\,18 & 0.036201\,@\,10 & 0.043447\,@\,7.5 & 0.053462\,@\,6.5 & 0.066007\,@\,6 & 0.082387\,@\,6 & 0.106399\,@\,6 & 0.064808\,@\,6.5 & 0.802545 \\
$(3,3.5]$ & 0.285715\,@\,3.5 & 0.031881\,@\,18.5 & 0.036444\,@\,11 & 0.043832\,@\,8.5 & 0.053696\,@\,7.5 & 0.066110\,@\,7 & 0.082895\,@\,6.5 & 0.107984\,@\,6.5 & 0.068058\,@\,7 & 0.791615 \\
$(3.5,4]$ & 0.272436\,@\,4 & 0.031917\,@\,19.5 & 0.036661\,@\,11.5 & 0.044100\,@\,9 & 0.054101\,@\,8 & 0.066753\,@\,7.5 & 0.083809\,@\,7.5 & 0.111171\,@\,7.5 & 0.072084\,@\,8 & 0.788032 \\
$(4,4.5]$ & 0.261905\,@\,4.5 & 0.031901\,@\,20.5 & 0.036806\,@\,12.5 & 0.044317\,@\,10 & 0.054390\,@\,8.5 & 0.067218\,@\,8 & 0.085273\,@\,8 & 0.114081\,@\,8.5 & 0.076659\,@\,9 & 0.787550 \\
$(4.5,5]$ & 0.253334\,@\,5 & 0.031849\,@\,21 & 0.036998\,@\,13 & 0.044677\,@\,10.5 & 0.054813\,@\,9.5 & 0.067954\,@\,9 & 0.086432\,@\,8.5 & 0.118217\,@\,9 & 0.081704\,@\,10 & 0.790978 \\
$(5,5.5]$ & 0.246213\,@\,5.5 & 0.031776\,@\,22 & 0.037131\,@\,13.5 & 0.044961\,@\,11 & 0.055338\,@\,10 & 0.068837\,@\,9.5 & 0.088316\,@\,9.5 & 0.121883\,@\,10 & 0.087168\,@\,11 & 0.796623 \\
\bottomrule
\end{tabular}
\end{center}
\end{landscape}

\begin{landscape}
\scriptsize
\setlength{\tabcolsep}{2.8pt}
\renewcommand{\arraystretch}{1.08}
\begin{center}
\textbf{Table D.2. Rational branch, $5.5<\alpha\le11$.}\par\medskip
\begin{tabular}{c|c|cccccccc|c}
\toprule
$\alpha$ slab & $E$ & $U_1$ & $U_2$ & $U_3$ & $U_4$ & $U_5$ & $U_6$ & $U_7$ & $U_8$ & total\\
\midrule
$(5.5,6]$ & 0.240197\,@\,6 & 0.031677\,@\,22.5 & 0.037220\,@\,14 & 0.045185\,@\,11.5 & 0.055770\,@\,10.5 & 0.069572\,@\,10 & 0.089974\,@\,10 & 0.126364\,@\,10.5 & 0.093017\,@\,12 & 0.803976 \\
$(6,6.5]$ & 0.235043\,@\,6.5 & 0.031563\,@\,23.5 & 0.037323\,@\,15 & 0.045394\,@\,12.5 & 0.056128\,@\,11 & 0.070211\,@\,11 & 0.091435\,@\,11 & 0.130327\,@\,11.5 & 0.099442\,@\,12.5 & 0.811866 \\
$(6.5,7]$ & 0.230577\,@\,7 & 0.031442\,@\,24 & 0.037418\,@\,15.5 & 0.045677\,@\,13 & 0.056544\,@\,12 & 0.071179\,@\,11.5 & 0.093388\,@\,11.5 & 0.135032\,@\,12 & 0.106187\,@\,13.5 & 0.822444 \\
$(7,7.5]$ & 0.226667\,@\,7.5 & 0.031304\,@\,24.5 & 0.037484\,@\,16 & 0.045916\,@\,13.5 & 0.057031\,@\,12.5 & 0.072026\,@\,12 & 0.095117\,@\,12 & 0.139270\,@\,12.5 & 0.113238\,@\,14.5 & 0.833053 \\
$(7.5,8]$ & 0.223215\,@\,8 & 0.031163\,@\,25.5 & 0.037529\,@\,16.5 & 0.046118\,@\,14 & 0.057457\,@\,13 & 0.072771\,@\,12.5 & 0.096656\,@\,12.5 & 0.143889\,@\,13.5 & 0.120591\,@\,15.5 & 0.844389 \\
$(8,8.5]$ & 0.220143\,@\,8.5 & 0.031021\,@\,26 & 0.037556\,@\,17 & 0.046290\,@\,14.5 & 0.057830\,@\,13.5 & 0.073432\,@\,13 & 0.098390\,@\,13.5 & 0.148335\,@\,14 & 0.128242\,@\,16.5 & 0.856239 \\
$(8.5,9]$ & 0.217392\,@\,9 & 0.030873\,@\,26.5 & 0.037595\,@\,18 & 0.046438\,@\,15 & 0.058160\,@\,14 & 0.074178\,@\,14 & 0.100132\,@\,14 & 0.152721\,@\,15 & 0.136253\,@\,17 & 0.868742 \\
$(9,9.5]$ & 0.214913\,@\,9.5 & 0.030719\,@\,27 & 0.037635\,@\,18.5 & 0.046597\,@\,16 & 0.058453\,@\,14.5 & 0.074981\,@\,14.5 & 0.101719\,@\,14.5 & 0.157314\,@\,15.5 & 0.144660\,@\,18 & 0.881991 \\
$(9.5,10]$ & 0.212667\,@\,10 & 0.030567\,@\,28 & 0.037662\,@\,19 & 0.046783\,@\,16.5 & 0.058824\,@\,15.5 & 0.075709\,@\,15 & 0.103168\,@\,15 & 0.161580\,@\,16 & 0.153343\,@\,19 & 0.895303 \\
$(10,10.5]$ & 0.210623\,@\,10.5 & 0.030420\,@\,28.5 & 0.037678\,@\,19.5 & 0.046947\,@\,17 & 0.059192\,@\,16 & 0.076373\,@\,15.5 & 0.104742\,@\,16 & 0.166080\,@\,17 & 0.162303\,@\,20 & 0.909358 \\
$(10.5,11]$ & 0.208755\,@\,11 & 0.030272\,@\,29 & 0.037685\,@\,20 & 0.047093\,@\,17.5 & 0.059527\,@\,16.5 & 0.076980\,@\,16 & 0.106338\,@\,16.5 & 0.170479\,@\,17.5 & 0.171539\,@\,21 & 0.923668 \\
\bottomrule
\end{tabular}
\end{center}
\end{landscape}

\begin{landscape}
\scriptsize
\setlength{\tabcolsep}{2.8pt}
\renewcommand{\arraystretch}{1.08}
\begin{center}
\textbf{Table D.3. Logarithmic branch, $11<\alpha\le17$.}\par\medskip
\begin{tabular}{c|c|cccccccc|c}
\toprule
$\alpha$ slab & $E$ & $U_1$ & $U_2$ & $U_3$ & $U_4$ & $U_5$ & $U_6$ & $U_7$ & $U_8$ & total\\
\midrule
$(11,11.5]$ & 0.203347\,@\,11.5 & 0.029354\,@\,29.5 & 0.036396\,@\,20.5 & 0.045134\,@\,18 & 0.056523\,@\,17 & 0.072357\,@\,16.5 & 0.098912\,@\,17 & 0.156956\,@\,18 & 0.156835\,@\,21.5 & 0.870814 \\
$(11.5,12]$ & 0.199912\,@\,12 & 0.028955\,@\,30 & 0.035964\,@\,21 & 0.044558\,@\,18.5 & 0.055712\,@\,17.5 & 0.071173\,@\,17 & 0.097365\,@\,17.5 & 0.154970\,@\,18.5 & 0.157446\,@\,22 & 0.861055 \\
$(12,12.5]$ & 0.196729\,@\,12.5 & 0.028581\,@\,30.5 & 0.035559\,@\,21.5 & 0.044019\,@\,19 & 0.054950\,@\,18 & 0.070058\,@\,17.5 & 0.095898\,@\,18 & 0.153232\,@\,19.5 & 0.158279\,@\,23 & 0.852305 \\
$(12.5,13]$ & 0.193767\,@\,13 & 0.028231\,@\,31 & 0.035180\,@\,22 & 0.043511\,@\,19.5 & 0.054230\,@\,18.5 & 0.069073\,@\,18.5 & 0.094502\,@\,18.5 & 0.151717\,@\,20 & 0.159062\,@\,24 & 0.844273 \\
$(13,13.5]$ & 0.191001\,@\,13.5 & 0.027901\,@\,31.5 & 0.034824\,@\,22.5 & 0.043032\,@\,20 & 0.053549\,@\,19 & 0.068170\,@\,19 & 0.093171\,@\,19 & 0.150235\,@\,20.5 & 0.160081\,@\,24.5 & 0.836964 \\
$(13.5,14]$ & 0.188412\,@\,14 & 0.027590\,@\,32 & 0.034487\,@\,23 & 0.042578\,@\,20.5 & 0.052903\,@\,19.5 & 0.067308\,@\,19.5 & 0.091898\,@\,20 & 0.148782\,@\,21 & 0.161043\,@\,25.5 & 0.830001 \\
$(14,14.5]$ & 0.185980\,@\,14.5 & 0.027296\,@\,32.5 & 0.034169\,@\,23.5 & 0.042148\,@\,21 & 0.052287\,@\,20 & 0.066484\,@\,20 & 0.090822\,@\,20.5 & 0.147357\,@\,21.5 & 0.162142\,@\,26 & 0.823685 \\
$(14.5,15]$ & 0.183690\,@\,15 & 0.027018\,@\,33 & 0.033867\,@\,24 & 0.041738\,@\,21.5 & 0.051700\,@\,20.5 & 0.065695\,@\,20.5 & 0.089782\,@\,21 & 0.146154\,@\,22.5 & 0.163237\,@\,27 & 0.817881 \\
$(15,15.5]$ & 0.181529\,@\,15.5 & 0.026753\,@\,33.5 & 0.033580\,@\,24.5 & 0.041347\,@\,22 & 0.051139\,@\,21 & 0.064937\,@\,21 & 0.088778\,@\,21.5 & 0.145024\,@\,23 & 0.164392\,@\,27.5 & 0.812479 \\
$(15.5,16]$ & 0.179485\,@\,16 & 0.026502\,@\,34 & 0.033309\,@\,25.5 & 0.040973\,@\,22.5 & 0.050601\,@\,21.5 & 0.064209\,@\,21.5 & 0.087805\,@\,22 & 0.143901\,@\,23.5 & 0.165585\,@\,28.5 & 0.807370 \\
$(16,16.5]$ & 0.177548\,@\,16.5 & 0.026263\,@\,34.5 & 0.033055\,@\,26 & 0.040615\,@\,23 & 0.050085\,@\,22 & 0.063507\,@\,22 & 0.086861\,@\,22.5 & 0.142787\,@\,24 & 0.166783\,@\,29 & 0.802504 \\
$(16.5,17]$ & 0.175708\,@\,17 & 0.026034\,@\,35 & 0.032812\,@\,26.5 & 0.040271\,@\,23.5 & 0.049621\,@\,23 & 0.062830\,@\,22.5 & 0.085946\,@\,23 & 0.141682\,@\,24.5 & 0.168050\,@\,30 & 0.797954 \\
\bottomrule
\end{tabular}
\end{center}
\end{landscape}

\section{Certification and reproducibility}\label{app:arithmetic}

The computer-assisted part of the proof is a fixed finite certificate, not a numerical optimization oracle.  Its proof object consists of the following files, archived with the arXiv source:
\begin{enumerate}[label=(\roman*),leftmargin=2em]
\item \path{Scalar_Certificate_HalfUnit.csv}, containing the $34$ slabs, the proposed maximizing half-integer for each of the $272$ nodal terms, an upward-rounded value bound, and the two directed ratio enclosures required by Corollary~\ref{cor:ratio};
\item \path{verify_scalar_certificate_exact.py}, a fixed-candidate checker using only Python integer arithmetic and \texttt{fractions.Fraction};
\item \path{verify_fixed_scalar_certificate.py}, an independent interval-arithmetic cross-check of the same fixed candidates;
\item \path{verify_T0_certificate.py} and \path{T0_Certificate.csv}, which verify the finite rational cover in Table~\ref{tab:T0};
\item \path{generate_scalar_certificate.py}, a candidate-generation program included for reproducibility.  It is not used by the fixed-candidate proof checker.
\end{enumerate}

The exact checker encloses logarithms by
\[
 \log\frac{1+z}{1-z}
 =2\sum_{m=0}^{N}\frac{z^{2m+1}}{2m+1}+R_N,
 \qquad
 |R_N|<\frac{2|z|^{2N+3}}{(2N+3)(1-z^2)}.
\]
For $11\le x\le17$, range reduction
\[
 \log x=4\log2+\log(x/16)
\]
keeps $|z|\le5/27$; for $\log2$ one has $z=1/3$.  Integer powers are exact.  A half-integer power of a positive rational interval is reduced to an integer power times a square root, and each square root is enclosed by rational decimal endpoints whose squares are checked exactly.  Thus the script verifies every ratio sign and every upper value using exact rational inequalities.

The smallest ratio margins in the fixed certificate are still comfortably resolved:
\[
 \min\bigl(R(s_\ast-1/2)-1\bigr)
 >4.7694\times10^{-6},
\]
where the minimum is taken only when a predecessor exists, and
\[
 \min\bigl(1-R(s_\ast)\bigr)
 >6.5842\times10^{-5}.
\]
The exact checker verifies all $272$ nodal maxima and reproduces the worst row
\[
 (0,1/2],\qquad F<0.966537.
\]
The interval checker provides a separate implementation of the same fixed-candidate verification.  The exact-rational $T_0$ checker verifies all $15$ rows in Table~\ref{tab:T0}, including both $g(U)<G$ and $C(L)G<3/200$.

For archive integrity, the source-package file \path{README_arXiv.txt} lists the SHA--256 hashes of the CSV certificate and all verification scripts.  These hashes identify the exact proof objects used to produce the bounds reported here.

The commands
\begin{verbatim}
python verify_scalar_certificate_exact.py
python verify_fixed_scalar_certificate.py
python verify_T0_certificate.py
\end{verbatim}
reproduce all proof-critical finite checks.  The generator may be run separately to reproduce the CSV; the proof checkers verify the fixed certificate independently of that generation step.

\section*{Code and data availability}
The exact scalar certificate, the fixed-candidate verifiers, the first-interval certificate, and the corresponding verification scripts are included in the source package of this arXiv version.  The versioned arXiv source package is the persistent archive of the proof objects used in this article.  The accompanying \path{README_arXiv.txt} records SHA--256 hashes for the certificate files and verification programs.  The candidate-generation program is included for reproducibility but is not used by either proof checker.

\section*{Acknowledgments}
The author thanks Lech Mazur for making the proof architecture and formal evidence publicly available, and Terence Tao for the streamlined scalar reduction and exposition that made the present optimization argument possible.

\section*{Declaration on the use of large language models}
Large language models played a substantial role in the development of this work.  OpenAI's ChatGPT, including GPT-5-series models used during 2026, was used to propose and compare proof architectures, derive candidate algebraic identities, generate exploratory symbolic and interval-arithmetic code, assist with selected Lean checks of auxiliary identities, and edit the manuscript.  The author formulated the research questions, selected and corrected the mathematical arguments, independently reran the computations, and takes full responsibility for the theorem, proof, numerical certificate, and exposition.  No language-model output is treated as mathematical evidence; the proof rests on the displayed analysis and the reproducible finite certificate.

\begin{thebibliography}{99}

\bibitem{BrownXiang1999}
J.~E. Brown and G.~Xiang,
\emph{Proof of the Sendov conjecture for polynomials of degree at most eight},
J. Math. Anal. Appl. \textbf{232} (1999), 272--292.

\bibitem{Chalebgwa2020}
T.~P. Chalebgwa,
\emph{Sendov's conjecture: a note on a paper of D\'egot},
Ann. Math. Sil. \textbf{35} (2021), 138--153.

\bibitem{Degot2014}
J.~D\'egot,
\emph{Sendov conjecture for high degree polynomials},
Proc. Amer. Math. Soc. \textbf{142} (2014), 1337--1349.

\bibitem{Marden1966}
M.~Marden,
\emph{Geometry of Polynomials}, 2nd ed.,
Mathematical Surveys No.~3, American Mathematical Society, Providence, 1966.

\bibitem{Mazur2026}
L.~Mazur,
\emph{A Computer-Assisted Proof of Sendov's Conjecture},
public version of 5 August 2026, ProofAtlas,
\url{https://proofatlas.ai/formalizations/sendov-conjecture/}.

\bibitem{MeirSharma1969}
A.~Meir and A.~Sharma,
\emph{On Ilyeff's conjecture},
Pacific J. Math. \textbf{31} (1969), 459--467.

\bibitem{RahmanSchmeisser2002}
Q.~I. Rahman and G.~Schmeisser,
\emph{Analytic Theory of Polynomials},
London Mathematical Society Monographs, Oxford University Press, 2002.

\bibitem{Tao2022}
T.~Tao,
\emph{Sendov's conjecture for sufficiently high degree polynomials},
Acta Math. \textbf{229} (2022), 347--392.

\bibitem{TaoDigestion2026}
T.~Tao,
\emph{A digestion of the proof of Sendov's conjecture},
What's New, 12 August 2026,
\url{https://terrytao.wordpress.com/2026/08/12/a-digestion-of-the-proof-of-sendovs-conjecture/}.

\bibitem{TaoLean2026}
T.~Tao,
\emph{Sendov: a Lean formalization of Sendov's conjecture},
GitHub repository, 2026,
\url{https://github.com/teorth/sendov}.

\end{thebibliography}

\end{document}
